\documentclass[12pt]{extarticle}

\usepackage{tabularx}
\usepackage{amsmath}
\usepackage{extsizes}
\usepackage{graphicx}
\usepackage{amssymb}
\usepackage[english,russian]{babel}
\usepackage[utf8]{inputenc}
\usepackage[a4paper]{geometry}
\geometry{top=20mm,bottom=30mm,left=20mm,right=15mm}
\usepackage[T1,T2A]{fontenc}
\usepackage{setspace}
\usepackage{titlesec}
\usepackage{indentfirst}
\usepackage{caption}
\renewcommand{\baselinestretch}{1.5}
\usepackage{spverbatim}
\usepackage{listings}
\usepackage{xcolor}
\usepackage{tabularray}
\lstset{extendedchars=\true}
\usepackage{hyperref}

\hypersetup{
    colorlinks=true,
    linkcolor=black,
    urlcolor=blue,
    pdftitle={Практическая работа №2: клиент Википедии},
}

\definecolor{sessionbg}{gray}{0.96}
\definecolor{sessionframe}{gray}{0.55}

\lstdefinestyle{session}{
    basicstyle=\ttfamily\small,
    breaklines=true,
    columns=fullflexible,
    keepspaces=true,
    frame=single,
    framesep=6pt,
    framerule=0.4pt,
    rulecolor=\color{sessionframe},
    backgroundcolor=\color{sessionbg},
    showstringspaces=false,
    upquote=true,
    aboveskip=8pt,
    belowskip=8pt,
    captionpos=b,
    xleftmargin=0pt,
    xrightmargin=0pt,
}

\newcommand{\cmd}[1]{\texttt{#1}}
\newcommand{\ui}[1]{\texttt{>\,#1}}

\definecolor{codebg}{gray}{0.96}
\definecolor{comment}{gray}{0.45}
\definecolor{keyword}{RGB}{0, 75, 150}
\definecolor{string}{RGB}{150, 80, 30}
\definecolor{number}{gray}{0.5}
\definecolor{identifier}{rgb}{0,0,0}

\lstdefinestyle{coolprint}{
  basicstyle   = \ttfamily\footnotesize,
  keywordstyle = \bfseries\color{keyword},
  commentstyle = \itshape\color{comment},
  stringstyle  = \color{string},
  identifierstyle = \color{identifier},
  backgroundcolor = \color{codebg},
  numbers      = left,
  numberstyle  = \footnotesize\color{number},
  numbersep    = 5pt,
  frame        = single,
  framerule    = 0.8pt,
  framesep     = 3pt,
  breaklines   = true,
  breakatwhitespace = true,
  tabsize      = 4,
  showspaces   = false,
  showstringspaces = false,
  showtabs     = false,
  xleftmargin  = 10pt,
  xrightmargin = 5pt,
  aboveskip    = \medskipamount,
  belowskip    = \medskipamount,
  captionpos   = b,
}

\usepackage{ragged2e}
\justifying

\titleformat{\section}{\centering\bfseries}{\thesection}{1em}{}
\titleformat{\subsection}{\centering\bfseries}{\thesubsection}{1em}{}

\begin{document}

\begin{center}
  \textbf{МИНОБРНАУКИ РОССИИ} \\
  \textbf{Санкт-Петербургский государственный} \\
  \textbf{электротехнический университет} \\
  \textbf{«ЛЭТИ» им. В.И. Ульянова (Ленина)} \\
  \textbf{Кафедра САПР} \\
\end{center}

\vspace{0.25\textheight}

\begin{center}
\textbf{Практическая работа №2} \\
\textbf{по дисциплине} \\
\textbf{«Объектно-ориентированное программирование»} \\
\end{center}

\vspace*{\fill}

\begin{center}
\begin{tabularx}{\textwidth}{ l X r }
 Студент гр. 5352 & \rule{8cm}{0.2mm} & Соболевский М.А. \\
 Преподаватель & \rule{8cm}{0.2mm} & Кулагин М.В. \\
\end{tabularx}
\end{center}

\vspace{3em}

\begin{center}
Санкт-Петербург\\
\the\year
\end{center}

\thispagestyle{empty}

\pagebreak

\section{Задание}

Разработать консольное приложение на языке~C, взаимодействующее с
внешним HTTP-API по протоколу HTTPS. Приложение должно:

\begin{itemize}
    \item реализовывать асинхронное (неблокирующее) взаимодействие с
          сетевыми ресурсами, не блокируя пользовательский интерфейс
          на время выполнения запроса;
    \item поддерживать интерактивный интерфейс пользователя
          (REPL-цикл), позволяющий отправлять запросы и получать
          результаты без перезапуска программы;
    \item корректно обрабатывать сетевые ошибки: недоступность
          сервера, ошибки TLS-рукопожатия, ошибки парсинга HTTP и
          JSON, таймауты;
    \item отображать прогресс выполнения длительных операций
          (индикатор активности, не засоряющий вывод);
    \item обеспечивать чистое разделение ответственности между
          слоями приложения: транспортом, клиентом API и
          пользовательским интерфейсом;
    \item демонстрировать применение принципов объектно-ориентированного
          проектирования в условиях языка~C: инкапсуляция состояния,
          композиция компонентов, полиморфизм через интерфейсы,
          абстракция.
\end{itemize}

В качестве целевого API выбран поисковый API Википедии: программа
принимает от пользователя поисковый запрос, выполняет HTTPS-запрос
к \texttt{ru.wikipedia.org}, отображает список найденных статей и
по выбору пользователя открывает соответствующую статью в браузере
системы.

\section{Спецификация программы}

Архитектура программы представлена на рисунке~\ref{fig:components}.
Приложение разделено на четыре слоя, каждый из которых знает только
о соседнем нижнем слое и не имеет обратных зависимостей.

\includegraphics[width=\textwidth]{assets/uml}

\begin{figure}[h]
\caption{Диаграмма компонентов}
\label{fig:components}
\end{figure}

Жизненный цикл асинхронного запроса показан на
рисунке~\ref{fig:request-fsm}. Каждый сетевой запрос выполняется в
отдельной корутине, которая приостанавливается на каждой
потенциально блокирующей операции и возобновляется event-loop'ом
при готовности сокета.

\includegraphics[width=\textwidth]{assets/wikipedia-cli-request-fsm}

\begin{figure}[h]
\caption{Состояния корутины сетевого запроса}
\label{fig:request-fsm}
\end{figure}

\subsection*{Слои приложения}

\begin{enumerate}
    \item \textbf{Транспортный слой} (\texttt{event\_loop2},
          \texttt{eloop\_net}, \texttt{eloop\_ssl},
          \texttt{eloop\_http})~--- внешняя библиотека, построенная
          на корутинах поверх \texttt{libco} и системного вызова
          \texttt{poll}. Обеспечивает неблокирующие сокеты,
          TLS-соединение (через OpenSSL) и HTTP-клиент с
          инкрементальным парсером заголовков. Слой не знает, как
          используется: приложение вызывает функции вида
          \texttt{http\_send\_req}, \texttt{http\_recv\_resp}, и
          получает данные по мере их поступления.

    \item \textbf{Уровень клиента API} (\texttt{wikipedia.c})~---
          один метод \texttt{request\_wikipedia}, который
          выполняет полный цикл: соединение, отправку
          HTTP-запроса, чтение тела ответа, парсинг JSON и возврат
          результата вызывающему. Всё взаимодействие с сетью
          инкапсулировано внутри; вызывающий видит только
          \texttt{struct wikipedia\_resp}.

    \item \textbf{Уровень пользовательского интерфейса}
          (\texttt{userloop.c}, \texttt{spinner.c})~---
          REPL-цикл, индикатор прогресса, отображение результатов
          и выбор пользователя. Не знает ни о сокетах, ни об
          HTTP, ни об SSL; работает только с
          \texttt{struct wikipedia\_resp}.

    \item \textbf{Вспомогательный слой} (\texttt{io/bufio.c},
          \texttt{strutils.c})~--- буферизация ввода-вывода и
          утилиты (URL-кодирование, отсечение пробелов).
\end{enumerate}

\subsection*{Асинхронность без колбэков}

Ключевое архитектурное решение~--- использование корутин вместо
колбэков. Каждый сетевой запрос оформляется как обычная функция,
внутри которой вызовы \texttt{esock\_*} и \texttt{http\_*}
выглядят как синхронные, но на самом деле приостанавливают
корутину и возвращают управление event-loop'у до готовности
сокета. Например, тело функции \texttt{wikipedia\_request}
читается сверху вниз, как обычный синхронный код, но при этом
работает параллельно с корутиной спиннера и корутиной основного
REPL-цикла.

Это позволяет писать асинхронный код без вложенных лямбд и без
явного управления состоянием парсера: корутина сама несёт в
своём стеке всё промежуточное состояние (соединение, буферы,
парсер JSON), а event-loop заботится только о том, чтобы
возобновить её, когда сокет готов.

\section{Описание интерфейса пользователя программы}

\lstset{style=session}

\subsection{Общие сведения}

Программа представляет собой интерактивный консольный клиент
поиска по Википедии. Взаимодействие построено по схеме REPL
(read--eval--print loop): программа выводит приглашение \texttt{>},
читает строку со стандартного ввода, выполняет поиск и снова
возвращается к приглашению.

Программа не принимает аргументов командной строки. Все действия
выполняются через интерактивный ввод. Завершение~--- по сигналу
EOF (комбинация \texttt{Ctrl+D}) или принудительно \texttt{Ctrl+C}.
Никаких файлов на диске программа не создаёт и не читает.

\subsection{Запуск и завершение}

Запуск:

\renewcommand{\baselinestretch}{1}
\begin{lstlisting}
$ ./wikipedia-cli
> _
\end{lstlisting}
\renewcommand{\baselinestretch}{1.5}

После запуска программа выводит приглашение и ожидает ввода
поискового запроса. Ввод завершается нажатием \texttt{Enter}.
Пустая строка игнорируется только на уровне поиска (Википедия
вернёт пустой результат).

Завершение работы~--- комбинацией \texttt{Ctrl+D} (EOF на
стандартном вводе). Программа корректно закрывает SSL-контекст,
освобождает память и выходит с кодом~0.

\subsection{Сценарий поиска}

Поиск выполняется в четыре шага.

\textbf{Шаг 1. Ввод запроса.} Пользователь вводит запрос после
приглашения:

\renewcommand{\baselinestretch}{1}
\begin{lstlisting}
> functional programming
\end{lstlisting}
\renewcommand{\baselinestretch}{1.5}

\textbf{Шаг 2. Выполнение запроса.} Пока идёт HTTPS-соединение и
сервер формирует ответ, программа отображает вращающийся индикатор
активности~--- символы \texttt{-}, \texttt{\textbackslash},
\texttt{|}, \texttt{/}, сменяющие друг друга с периодом
100~мс. Индикатор рисуется в одной позиции строки с помощью
символа возврата каретки \texttt{\textbackslash r}, поэтому не
засоряет вывод:

\renewcommand{\baselinestretch}{1}
\begin{lstlisting}
> functional programming
Searching... -
\end{lstlisting}
\renewcommand{\baselinestretch}{1.5}

\textbf{Шаг 3. Вывод результатов.} По завершении запроса
индикатор стирается, и программа выводит пронумерованный список
найденных статей:

\renewcommand{\baselinestretch}{1}
\begin{lstlisting}
1: Функциональное программирование
2: Лямбда-исчисление
3: Haskell
4: OCaml
5: Функтор
6: Монада (программирование)
7: Чистота (программирование)
8: Ленивые вычисления
9: Комбинаторная логика
10: Схема (язык программирования)
\end{lstlisting}
\renewcommand{\baselinestretch}{1.5}

Википедия по умолчанию возвращает 10 результатов на запрос.

\textbf{Шаг 4. Выбор страницы.} Программа выводит приглашение
\texttt{select page>} и ожидает ввода номера из списка:

\renewcommand{\baselinestretch}{1}
\begin{lstlisting}
select page> 3
\end{lstlisting}
\renewcommand{\baselinestretch}{1.5}

После ввода корректного номера программа запускает браузер
(\texttt{xdg-open}) на соответствующей странице Википедии и
возвращается к основному приглашению:

\renewcommand{\baselinestretch}{1}
\begin{lstlisting}
select page> 3
> _
\end{lstlisting}
\renewcommand{\baselinestretch}{1.5}

Если пользователь введёт \texttt{q} вместо номера, выбор будет
отменён, и программа вернётся к основному приглашению без
открытия браузера.

\subsection{Спиннер и его роль}

Спиннер~--- отдельная корутина, которая запускается на время
сетевого запроса и останавливается, когда приходит ответ. С точки
зрения пользователя это означает, что интерфейс остаётся
отзывчивым во время долгого соединения: каждая итерация спиннера
происходит примерно раз в 100~мс, тогда как основной REPL-цикл в
это время заблокирован на ожидании сетевого ответа (но не
процессор~--- операционная система при этом свободна).

\subsection{Обработка ошибочного ввода}

\textbf{Нечисловой номер страницы.} Если в ответ на
\texttt{select page>} ввести строку, которая не начинается с
цифры, преобразование \texttt{strtol} вернёт~0, что меньше 1 и
поэтому не попадёт в допустимый диапазон:

\renewcommand{\baselinestretch}{1}
\begin{lstlisting}
select page> abc
entered number is not in the list
select page> _
\end{lstlisting}
\renewcommand{\baselinestretch}{1.5}

\textbf{Номер вне диапазона.} Аналогично, число больше количества
найденных статей или меньше единицы отклоняется с тем же
сообщением.

\textbf{Отмена выбора.} Ввод \texttt{q} на приглашении
\texttt{select page>} трактуется как отказ от открытия страницы.
Программа возвращается в главное приглашение.

\subsection{Сетевые ошибки}

Если на этапе соединения или TLS-рукопожатия произошла ошибка,
программа печатает диагностику в стандартный поток ошибок и
возвращается к приглашению:

\renewcommand{\baselinestretch}{1}
\begin{lstlisting}
> test
Searching... error: connection failed
> _
\end{lstlisting}
\renewcommand{\baselinestretch}{1.5}

Различаются три класса ошибок:

\begin{itemize}
    \item \texttt{error: connection failed}~--- не удалось
          установить TCP-соединение с сервером (DNS не разрешился,
          сервер недоступен, порт закрыт);
    \item \texttt{error: transport failure}~--- соединение
          установлено, но произошла ошибка на уровне TLS или
          HTTP (обрыв соединения, некорректный ответ сервера);
    \item \texttt{error: invalid response}~--- сервер вернул
          HTTP-код, отличный от 200 (например, 4xx или 5xx).
\end{itemize}

\subsection{Модель данных}

Результат одного запроса описывается структурой
\texttt{struct wikipedia\_resp}, которая содержит:

\begin{itemize}
    \item \texttt{status}~--- перечисление
          \texttt{enum wiki\_resp\_status}, одно из
          \texttt{WRS\_OK}, \texttt{WRS\_CONNECTION\_FAILED},
          \texttt{WRS\_INVALID\_RESPONSE},
          \texttt{WRS\_TRANSPORT\_FAILURE};
    \item \texttt{payload}~--- указатель на разобранный JSON-ответ
          (тип \texttt{json\_object*}), действителен только при
          \texttt{status == WRS\_OK}.
\end{itemize}

Вызывающий код обязан проверять \texttt{status} перед обращением
к \texttt{payload}. Это часть неявного контракта между
\texttt{request\_wikipedia} и её вызывающим.

\section{Примеры работы и описание}

Ниже приведены типовые сценарии работы. Строки, начинающиеся с
\texttt{>\ },~--- это то, что вводит пользователь; остальные
строки~--- вывод программы.

\subsection{Успешный поиск и открытие страницы}

\renewcommand{\baselinestretch}{1}
\begin{lstlisting}
$ ./wikipedia-cli
> корутина
Searching... \
1: Сопрограмма
2: Корутина
3: Асинхронный ввод-вывод
4: Fibers
5: Поток выполнения
6: Зелёные потоки
7: Легковесные процессы
8: Concurrency
9: Setcontext
10: Ucontext
select page> 1
> _
\end{lstlisting}
\renewcommand{\baselinestretch}{1.5}

После выбора страницы браузер открывается на
\texttt{ru.wikipedia.org/w/index.php?curid=XXXXXXX}. Сообщений об
успехе программа не печатает~--- открытие браузера считается
само собой разумеющимся.

\subsection{Отмена выбора страницы}

\renewcommand{\baselinestretch}{1}
\begin{lstlisting}
> haskell
Searching... |
1: Haskell
2: GHC
3: Glasgow Haskell Compiler
...
select page> q
> _
\end{lstlisting}
\renewcommand{\baselinestretch}{1.5}

Ввод \texttt{q} не открывает браузер и не печатает сообщений.

\subsection{Ошибка DNS или соединения}

Если Wi-Fi отключён и \texttt{ru.wikipedia.org} не резолвится:

\renewcommand{\baselinestretch}{1}
\begin{lstlisting}
> test
Searching... error: connection failed
> _
\end{lstlisting}
\renewcommand{\baselinestretch}{1.5}

Программа не завершается и не «зависает»~--- она возвращает
управление в REPL, и следующий запрос можно отправить сразу.

\subsection{Обрыв соединения во время запроса}

Если соединение было установлено, но сервер закрыл его до
завершения ответа:

\renewcommand{\baselinestretch}{1}
\begin{lstlisting}
> test
Searching... error: transport failure
> _
\end{lstlisting}
\renewcommand{\baselinestretch}{1.5}

Этот класс ошибок покрывает также ошибки TLS (например,
несовпадение сертификата) и ошибки парсинга HTTP-ответа.

\subsection{Сервер вернул 4xx/5xx}

Если сервер вернул HTTP-код, отличный от 200:

\renewcommand{\baselinestretch}{1}
\begin{lstlisting}
> test
Searching... error: invalid response
> _
\end{lstlisting}
\renewcommand{\baselinestretch}{1.5}

\subsection{Некорректный номер страницы}

\renewcommand{\baselinestretch}{1}
\begin{lstlisting}
> linux
Searching... /
1: Linux
2: Ядро Linux
3: GNU
...
select page> 42
entered number is not in the list
select page> abc
entered number is not in the list
select page> 2
> _
\end{lstlisting}
\renewcommand{\baselinestretch}{1.5}

\subsection{Завершение через Ctrl+D}

\renewcommand{\baselinestretch}{1}
\begin{lstlisting}
> quit
$
\end{lstlisting}
\renewcommand{\baselinestretch}{1.5}

Программа закрывает SSL-контекст, освобождает все выделенные
ресурсы и выходит с кодом~0. Слово \texttt{quit} здесь не
команда~--- программа получает \texttt{EOF} на стандартном вводе,
когда пользователь нажимает \texttt{Ctrl+D}; строка \texttt{quit}
в примере~--- это просто демонстрация того, что при этом
происходит.

\section{Текст программы}
\lstset{style=coolprint}
\renewcommand{\baselinestretch}{1}
\input{assets/source}
\renewcommand{\baselinestretch}{1.5}

\section{Выводы}

В ходе выполнения практической работы было разработано консольное
приложение~--- клиент поиска по Википедии~--- на языке~C~99 с
использованием собственной библиотеки асинхронного
ввода-вывода. Программа выполняет HTTPS-запросы, парсит JSON,
отображает результаты и открывает выбранную статью в браузере,
не блокируя пользовательский интерфейс.

\subsection*{Реализованная функциональность}

В соответствии с заданием были реализованы:

\begin{itemize}
    \item интерактивный REPL-цикл с обработкой \texttt{Ctrl+D}
          (EOF) как команды завершения;
    \item асинхронное HTTPS-соединение с
          \texttt{ru.wikipedia.org}, не блокирующее основной
          цикл приложения;
    \item парсинг HTTP-ответа (заголовки, тело, chunked
          encoding) и парсинг JSON-ответа через библиотеку
          \texttt{json-c};
    \item индикатор прогресса (спиннер) в виде отдельной
          корутины, обновляющийся каждые 100~мс;
    \item URL-кодирование поискового запроса;
    \item обработка ошибок трёх классов: сетевые, транспортные,
          ошибки валидации ответа;
    \item открытие браузера на выбранной статье через
          \texttt{xdg-open}.
\end{itemize}

\subsection*{Соблюдение принципов объектно-ориентированного
программирования}

\begin{enumerate}
    \item \textbf{Инкапсуляция.} Все структуры данных
          (\texttt{struct wikipedia\_resp}, \texttt{struct
          enet\_socket}, \texttt{struct bufio},
          \texttt{struct http\_transport}) объявлены как
          непрозрачные типы в заголовочных файлах; их поля
          видны только внутри реализующего модуля. Изменение
          состояния выполняется исключительно через
          функции-методы (\texttt{http\_send\_req},
          \texttt{esock\_recv}, \texttt{bio\_flush} и~т.~д.),
          которые проверяют предусловия и возвращают
          коды ошибок.

    \item \textbf{Композиция.} Приложение построено снизу
          вверх: \texttt{event\_loop} управляет корутинами;
          \texttt{eloop\_net} обеспечивает сокеты;
          \texttt{eloop\_ssl} оборачивает их в TLS;
          \texttt{eloop\_http} строит HTTP-клиент над
          транспортом через \texttt{bufio}; \texttt{wikipedia}
          инкапсулирует бизнес-логику запроса;
          \texttt{userloop} управляет интерфейсом. Каждый
          слой владеет своими подобъектами и освобождает их
          при завершении.

    \item \textbf{Полиморфизм.} В коде активно используются
          указатели на функции как способ абстракции над
          поведением:
          \begin{itemize}
              \item \texttt{bio\_init} принимает указатели на
                    \texttt{read} и \texttt{write}, что
                    позволяет HTTP-клиенту одинаково работать
                    поверх TCP и поверх TLS без изменения
                    своего кода;
              \item \texttt{esock\_setup\_transport} принимает
                    указатель на функцию очистки, что позволяет
                    подключать произвольные транспорты (TLS,
                    SOCKS5-прокси) к одному сокету;
              \item \texttt{esock\_callbacks} и
                    \texttt{essl\_callbacks} реализуют паттерн
                    «наблюдатель» для событий закрытия,
                    ошибок и таймаутов;
              \item \texttt{eloop\_add} принимает
                    \texttt{ev\_routine}~--- функцию без
                    аргументов,~--- что позволяет запускать
                    произвольный код как корутину.
          \end{itemize}
          Во всех случаях выбор конкретной реализации
          происходит в момент создания объекта, а не в момент
          вызова~--- классический паттерн «стратегия».

    \item \textbf{Абстракция.} Вызывающий код
          (\texttt{userloop.c}) не знает, что находится под
          \texttt{request\_wikipedia}: сокеты, TLS, HTTP,
          JSON~--- всё скрыто за одной функцией, которая
          принимает строку и возвращает структуру с
          результатом. Аналогично, HTTP-клиент не знает,
          работает он поверх TCP или TLS,~--- для него есть
          только \texttt{read}/\texttt{write} callbacks.
          Каждый слой скрывает детали нижнего слоя за
          минимальным интерфейсом.

    \item \textbf{Разделение ответственности.} Три уровня
          (транспорт, клиент API, интерфейс пользователя)
          ничего не знают друг о друге, кроме
          непосредственного соседа снизу. Изменение,
          например, формата вывода результатов не затрагивает
          ни HTTP-клиент, ни сетевой слой; добавление нового
          API-эндпоинта требует изменения только
          \texttt{wikipedia.c}.
\end{enumerate}

\subsection*{Принятые архитектурные решения}

\begin{itemize}
    \item \textbf{Корутины вместо колбэков.} Асинхронный код
          выглядит как обычный синхронный: нет вложенных
          лямбд, нет ручного state machine, нет явного
          управления контекстом парсера. Всё промежуточное
          состояние (соединение, буферы, парсер) лежит на
          стеке корутины и автоматически освобождается при её
          завершении.

    \item \textbf{Стек-выделенные замыкания.} Структура
          \texttt{struct wikipedia\_req\_closure} находится
          на стеке функции \texttt{request\_wikipedia}, но
          это безопасно: корутина приостанавливает сама себя
          на \texttt{ev\_recv\_message}, и её стек остаётся
          живым до получения ответа. Такая схема невозможна
          для потоков, но естественна для корутин.

    \item \textbf{setjmp/longjmp для обработки ошибок HTTP.}
          Вместо возврата кодов ошибок из каждой функции
          HTTP-слоя используется \texttt{longjmp}: если парсер
          обнаруживает ошибку, он «выбрасывает» её в
          \texttt{jmp\_buf}, который находится на стеке
          вызывающего кода. Это разделяет нормальный поток и
          обработку ошибок, но требует ручной очистки
          ресурсов~--- поэтому переменные \texttt{buf} и
          \texttt{jtk} объявлены \texttt{volatile} (после
          \texttt{longjmp} компилятор не обязан сохранять
          значения не-\texttt{volatile} локальных переменных).

    \item \textbf{Отдельный таск для спиннера.} Спиннер
          реализован не как «фоновый поток», а как корутина,
          которая крутится в том же event-loop. Пока
          HTTP-запрос ждёт данных от сокета, спиннер получает
          управление раз в 100~мс и перерисовывает символ.
          Это не блокирует event-loop и не требует
          дополнительных потоков или сигналов.

    \item \textbf{Транспортная абстракция HTTP.} Через
          \texttt{bio\_init} HTTP-клиент получает указатели на
          функции чтения и записи, которые могут указывать на
          сокет или на TLS-обёртку. Один и тот же парсер
          работает и для \texttt{http://}, и для
          \texttt{https://}, без \texttt{if}-ов по типу
          транспорта.
\end{itemize}

\subsection*{Полученные навыки}

\begin{itemize}
    \item проектирование многослойной архитектуры и
          распределение ответственности между модулями;
    \item реализация асинхронного ввода-вывода на основе
          корутин и неблокирующих сокетов;
    \item работа с протоколом TLS через OpenSSL;
    \item написание HTTP/1.1 клиента с инкрементальным
          парсером заголовков и поддержкой chunked encoding;
    \item парсинг JSON и обработка древовидных структур через
          \texttt{json-c};
    \item организация управления памятью при отсутствии
          сборщика мусора;
    \item применение паттерна «стратегия» и «наблюдатель» в
          условиях отсутствия классов;
    \item построение интерфейса пользователя с
          неблокирующим индикатором прогресса;
    \item работа с системой контроля версий Git и оформление
          проекта в виде отдельного репозитория.
\end{itemize}

\subsection*{Заключение}

Поставленная задача выполнена в полном объёме. Программа
функционирует согласно требованиям задания, поддерживает
интерактивный интерфейс, корректно обрабатывает сетевые и
пользовательские ошибки, демонстрирует применение принципов
объектно-ориентированного проектирования в проекте на языке~C.

Основное отличие от первой практики~--- работа с внешним
состоянием (сетью, таймаутами, ошибками ввода-вывода) и
необходимость поддерживать отзывчивость интерфейса при
длительных операциях. Это потребовало построения собственной
библиотеки асинхронного ввода-вывода на корутинах, что выходит
за рамки типичного учебного задания, но позволяет глубже
понять, как устроены современные async-фреймворки (libuv,
tokio, async/await) и как они соотносятся с классическими
принципами ООП.

Полученный код легко расширяется: добавление нового
API-эндпоинта сводится к написанию новой функции в
\texttt{wikipedia.c}; замена \texttt{xdg-open} на
платформо-специфичный запуск браузера~--- к правке одной
функции в \texttt{userloop.c}; замена формата вывода
результатов не затрагивает ни HTTP-клиент, ни сетевой слой.

\end{document}
